% theme: addition_theorems
\MajorQuestionLesson{加法定理}
\SetRowSpace{4mm}

\TypeQuestion{次の値を求めなさい。}{addition_theorems_exact_value}
\QQRowThree{$\sin75^\circ$}{}{$\tan15^\circ$}{}{$\cos105^\circ$}{}
\QQRowThree{$\cos\dfrac{\pi}{12}$}{}{$\sin\dfrac{7\pi}{12}$}{}{$\tan\dfrac{5\pi}{12}$}{}
\QQRowThree{$\tan105^\circ$}{}{$\sin(-15^\circ)$}{}{$\cos165^\circ$}{}
\QQRowThree{$\sin\dfrac{13\pi}{12}$}{}{$\cos\left(-\dfrac{5\pi}{12}\right)$}{}{$\tan\dfrac{11\pi}{12}$}{}

\TypeQuestion{次の式の値を求めなさい。}{addition_theorems_recognize}
\QQ{$\sin38^\circ\cos22^\circ+\cos38^\circ\sin22^\circ$}{}
\QQ{$\cos\dfrac{7\pi}{18}\cos\dfrac{\pi}{18}+\sin\dfrac{7\pi}{18}\sin\dfrac{\pi}{18}$}{}
\QQ{$\dfrac{\tan47^\circ-\tan17^\circ}{1+\tan47^\circ\tan17^\circ}$}{}
\QQ{$\cos83^\circ\cos37^\circ-\sin83^\circ\sin37^\circ$}{}
\QQ{$\sin\dfrac{5\pi}{18}\cos\dfrac{4\pi}{9}-\cos\dfrac{5\pi}{18}\sin\dfrac{4\pi}{9}$}{}
\QQ{$\dfrac{\tan13^\circ+\tan32^\circ}{1-\tan13^\circ\tan32^\circ}$}{}

\TypeQuestion{次の式を、$\sin\theta,\ \cos\theta$ を用いて表しなさい。}{addition_theorems_expand}
\QQRow{$\sin\left(\theta+\dfrac{\pi}{3}\right)$}{}{$\cos\left(\theta-\dfrac{\pi}{4}\right)$}{}
\QQRow{$\cos\left(\dfrac{\pi}{6}+\theta\right)$}{}{$\sin\left(\dfrac{\pi}{4}-\theta\right)$}{}
\QQRow{$\sin\left(\theta-\dfrac{2\pi}{3}\right)$}{}{$\cos\left(\theta+\dfrac{5\pi}{6}\right)$}{}

\LessonQuestionBreak

\TypeQuestion{次の条件のもとで、指定された式の値を求めなさい。}{addition_theorems_one_angle}
\QQ{$\sin\theta=\dfrac{3}{5}$、$\dfrac{\pi}{2}<\theta<\pi$ のとき、$\sin\left(\theta+\dfrac{\pi}{6}\right)$}{}
\QQ{$\cos\theta=\dfrac{5}{13}$、$\dfrac{3\pi}{2}<\theta<2\pi$ のとき、$\cos\left(\theta-\dfrac{\pi}{3}\right)$}{}
\QQ{$\tan\theta=2$、$0<\theta<\dfrac{\pi}{2}$ のとき、$\tan\left(\theta-\dfrac{\pi}{4}\right)$}{}
\QQ{$\cos\theta=-\dfrac{2}{3}$、$\pi<\theta<\dfrac{3\pi}{2}$ のとき、$\sin\left(\dfrac{\pi}{4}-\theta\right)$}{}
\QQ{$\tan\theta=-\dfrac{1}{2}$、$\dfrac{\pi}{2}<\theta<\pi$ のとき、$\cos\left(\theta+\dfrac{\pi}{4}\right)$}{}
\QQ{$\sin\theta=-\dfrac{5}{13}$、$\pi<\theta<\dfrac{3\pi}{2}$ のとき、$\tan\left(\theta+\dfrac{\pi}{6}\right)$}{}

\TypeQuestion{次の条件のもとで、指定された式の値を求めなさい。}{addition_theorems_two_angles}
\QQ{$\sin\alpha=\dfrac{3}{5}$、$\cos\beta=\dfrac{12}{13}$、$0<\alpha,\beta<\dfrac{\pi}{2}$ のとき、$\sin(\alpha+\beta)$}{}
\QQ{$\cos\alpha=-\dfrac{4}{5}$、$\sin\beta=\dfrac{5}{13}$、$\dfrac{\pi}{2}<\alpha<\pi$、$0<\beta<\dfrac{\pi}{2}$ のとき、$\cos(\alpha-\beta)$}{}
\QQ{$\tan\alpha=\dfrac{1}{3}$、$\tan\beta=-\dfrac{2}{5}$ のとき、$\tan(\alpha+\beta)$}{}
\QQ{$\sin\alpha=-\dfrac{4}{5}$、$\cos\beta=-\dfrac{5}{13}$、$\dfrac{3\pi}{2}<\alpha<2\pi$、$\pi<\beta<\dfrac{3\pi}{2}$ のとき、$\sin(\alpha-\beta)$}{}
\QQ{$\tan\alpha=-\dfrac{3}{4}$、$\tan\beta=\dfrac{1}{2}$ のとき、$\tan(\alpha-\beta)$}{}
\QQ{$\cos\alpha=\dfrac{1}{3}$、$\sin\beta=-\dfrac{2}{3}$、$0<\alpha<\dfrac{\pi}{2}$、$\dfrac{3\pi}{2}<\beta<2\pi$ のとき、$\cos(\alpha+\beta)$}{}

\TypeQuestion{$0<\alpha,\beta<\dfrac{\pi}{2}$ とする。次の条件のもとで、$\alpha+\beta$ の値を求めなさい。}{addition_theorems_determine_angle}
\QQRow{$\tan\alpha=2,\quad\tan\beta=3$}{}{$\tan\alpha=\dfrac{1}{2},\quad\tan\beta=2$}{}
\QQ{$\tan\alpha=\dfrac{1}{3},\quad\tan\beta=\dfrac{1}{2}$}{}

\LessonQuestionBreak

\TypeQuestion{次の等式を証明しなさい。ただし、各式に含まれる分母は $0$ でないものとする。}{addition_theorems_identity}
\QQ{$\sin(\alpha+\beta)\sin(\alpha-\beta)=\sin^2\alpha-\sin^2\beta$}{}
\QQ{$\cos(\alpha+\beta)\cos(\alpha-\beta)=\cos^2\alpha-\sin^2\beta$}{}
\QQ{$\tan(\alpha+\beta)-\tan\alpha=\dfrac{\sin\beta}{\cos(\alpha+\beta)\cos\alpha}$}{}
\QQ{$\sin(\alpha+\beta)\cos\beta-\cos(\alpha+\beta)\sin\beta=\sin\alpha$}{}

\TypeQuestion{次の方程式を、$0\leqq\theta<2\pi$ の範囲で解きなさい。}{addition_theorems_equation}
\QQRow{$\sin\left(\theta+\dfrac{\pi}{3}\right)=\sin\theta$}{}{$\cos\left(\theta-\dfrac{\pi}{6}\right)=\cos\theta$}{}
\QQRow{$\sin\left(\theta+\dfrac{\pi}{6}\right)=\cos\theta$}{}{$\cos\left(\theta+\dfrac{\pi}{3}\right)=\sin\theta$}{}
\QQ{$\cos\left(\theta+\dfrac{\pi}{4}\right)+\cos\left(\theta-\dfrac{\pi}{4}\right)=1$}{}
